% Paper Skeleton — Annotated LaTeX structure for academic papers
% Usage: Copy this skeleton and fill in each section following the paper-writing skill guidelines.

\documentclass{article}
% Adapt the document class to your venue (e.g., \documentclass{cvpr}, \documentclass[review]{acl})

\usepackage{booktabs}    % Professional table rules
\usepackage{graphicx}    % Figures
\usepackage{amsmath}     % Math
\usepackage{hyperref}    % Hyperlinks
\usepackage{xcolor}      % Color for highlights

% Color macros for table highlights
\definecolor{bestcolor}{HTML}{FF0000}
\definecolor{secondcolor}{HTML}{0000FF}
\newcommand{\best}[1]{\textcolor{bestcolor}{\textbf{#1}}}
\newcommand{\second}[1]{\textcolor{secondcolor}{\underline{#1}}}

\title{Your Paper Title}
\author{Author Names}

\begin{document}
\maketitle

% ==============================================================================
% ABSTRACT (Step 9)
% Answer: (1) Technical problem (2) Contribution (3) Why it works (4) Advantage
% Templates: abstract-templates.md (Version 1/2/3)
% ==============================================================================
\begin{abstract}
% Task. Technical challenge.
% [Choose Version 1, 2, or 3 from abstract-templates.md]
% Insight (if Version 2).
% Technical contribution.
% Benefits. Experiment summary.
\end{abstract}

% ==============================================================================
% INTRODUCTION (Step 2)
% Structure: Task → Previous methods → Challenge → Our pipeline → Advantages
% Templates: introduction-templates.md
% ==============================================================================
\section{Introduction}
\label{sec:intro}

% Part 1: Introduce task and application
% [Choose Version 1-4 from introduction-templates.md, Part 1]

% Part 2: Previous methods and technical challenge
% [Choose Version 1-3 from introduction-templates.md, Part 2]

% Part 3: Our pipeline and contributions
% [Choose Version 1-4 from introduction-templates.md, Part 3]
% Anti-pattern: Do NOT write "naive solution → our improvement"

% Contribution paragraph (typically bulleted or numbered):
% Our contributions are summarized as follows:
% \begin{itemize}
%   \item We propose ... which ...
%   \item We introduce ... that enables ...
%   \item Extensive experiments demonstrate ...
% \end{itemize}

% ==============================================================================
% RELATED WORK (Step 7)
% Process: List papers → Group into topics → Write paragraphs
% Guide: related-work-guide.md
% ==============================================================================
\section{Related Work}
\label{sec:related}

% \paragraph{Topic 1.} Survey + position our work.
% \paragraph{Topic 2.} Survey + position our work.
% \paragraph{Topic 3.} Survey + position our work.

% ==============================================================================
% METHOD (Step 3)
% Each module: Motivation → Module Design → Technical Advantages
% Templates: method-templates.md
% ==============================================================================
\section{Method}
\label{sec:method}

% Overview paragraph:
% Setting (1-2 sentences) → Core contribution (1-2 sentences) → Pipeline figure reference → Section roadmap

\subsection{Module 1}
\label{sec:module1}
% Motivation paragraph
% Module design paragraph(s): structure + forward process
% Technical advantages paragraph

\subsection{Module 2}
\label{sec:module2}
% Motivation paragraph
% Module design paragraph(s): structure + forward process
% Technical advantages paragraph

\subsection{Module 3}
\label{sec:module3}
% Motivation paragraph
% Module design paragraph(s): structure + forward process
% Technical advantages paragraph

% ==============================================================================
% EXPERIMENTS (Step 5)
% Answer: (1) Better than existing? (2) Modules effective? (3) Upper limit?
% Guide: experiments-guide.md
% ==============================================================================
\section{Experiments}
\label{sec:exp}

\subsection{Experimental Setup}
% Datasets, evaluation metrics, implementation details, baselines

\subsection{Comparison with State-of-the-Art Methods}
% Quantitative table + qualitative figure + analysis

\subsection{Ablation Study}
% Part 1: Core contribution ablation (one big table)
% Part 2: Module design choices (several small tables)

\subsection{Applications}  % Optional but recommended
% Demos on challenging data

% ==============================================================================
% CONCLUSION (Step 11)
% Must include Limitation — reviewers flag its absence
% ==============================================================================
\section{Conclusion}
\label{sec:conclusion}

% Summary of contributions and results.

% Limitation: task/setting limitations (NOT technical defects).
% Rule: "If not below SOTA metrics, it is not a technical defect."

\end{document}
